\documentclass[11pt]{article}
\usepackage[T1]{fontenc}
\usepackage{lmodern,amsmath,amssymb,amsthm,mathtools,microtype,booktabs}
\usepackage[margin=1in]{geometry}
\usepackage[hidelinks]{hyperref}
\usepackage{needspace}
\setlength{\emergencystretch}{2em}
\newtheorem{theorem}{Theorem}[section]
\newtheorem{lemma}[theorem]{Lemma}
\newtheorem{proposition}[theorem]{Proposition}
\newtheorem{corollary}[theorem]{Corollary}
\theoremstyle{remark}\newtheorem{remark}[theorem]{Remark}
\newcommand{\R}{\mathbb R}\newcommand{\Var}{\operatorname{Var}}
\newcommand{\Ptwo}{\mathcal P_2}\newcommand{\TV}{\operatorname{TV}}
\newcommand{\norm}[1]{\lVert#1\rVert}\newcommand{\ind}{\mathbf 1}
\title{Tail-sensitive convex-gradient interpolation\\and Brenier stability beyond bounded targets}
\author{mxym\thanks{Repository-account attribution: \url{https://github.com/mxym/math}. Prepared with AI assistance. No institutional affiliation is asserted.}}
\date{October 7, 2026\quad Version 2}
\hypersetup{pdfauthor={mxym},pdflang={en-US}}
\newcommand{\E}{\mathbb E}
\begin{document}\maketitle
\begin{abstract}
We prove a convex-gradient interpolation inequality with an explicit second-moment tail remainder, for densities of finite integrated coordinate-slice variation. It requires no global bound on the gradients. Under a bounded $p$th gradient moment, $p>2$, the resulting potential-to-gradient exponent is $(p-2)/(3p-2)$, and this interpolation exponent is sharp. Combining the inequality with a radius-independent potential stability estimate yields quantitative Brenier-map stability for unbounded targets with any $p>2$, independently of the ambient dimension as a moment threshold. The application to strongly log-concave sources uses the potential estimate in manuscript 001; compact log-concave sources are also covered by completion of its bounded-target estimate. A tail-infimum bound treats general uniformly integrable second moments, and stretched-exponential target tails give one-third stability with a logarithmic factor. A separate two-atom construction proves the general finite-moment transport exponent sharp for one fixed smooth, full-support, strongly log-concave source. In contrast, superquadratic product sources retain a sharp one-third exponent. Thus full support and smoothness alone do not replace quantitative translation regularity.
\end{abstract}
\paragraph{Status and dependencies.}
The interpolation theorem and its sharp example have complete proofs below, independent of the new transport results in the repository. The Brenier applications explicitly use the radius-independent potential estimate from manuscript 001~\cite{M001}; that proof has not been machine-formalized or independently peer reviewed. The bounded-gradient finite-difference mechanism is inherited from OpenAI family 374~\cite{O374} and 001. The additional argument restricts each coordinate slice to the interval where both derivatives are bounded and charges the complement to the actual gradient tails. Earlier potential-to-gradient approaches and unbounded-target results include Delalande--M\'erigot~\cite{DM}. A parallel update, manuscript 001 v3, was found during release reconciliation. It already contains the general log-concave target-tail bound, sharp truncated-source examples, and Gaussian and bounded-Hessian one-third estimates. Those overlapping transport bounds are not claimed as a distinct new result here. The present additional statements are the BV-density interpolation theorem, superquadratic-source translation estimates, and the fixed smooth full-support counterexample in Section~\ref{sec:steep}. No priority or journal-tier certification is asserted.

\section{A tail-sensitive interpolation theorem}
Let $D\subset\R^d$ be open and convex. Let $r\ge0$ be a probability density, zero almost everywhere outside $D$. For almost every $z\in e_j^\perp$, assume the coordinate slice $t\mapsto r(z+te_j)$ has a representative of bounded variation. Define
\begin{equation}\label{eq:slice-variation}
 V(r)=\frac12\sum_{j=1}^d\int_{e_j^\perp}
                  \TV_{\R}(r(z+\cdot e_j))\,dz<\infty.
\end{equation}
For a $BV$ density this is half the sum of its directional derivative total variations. For a log-concave density with finite slice functional it equals
\begin{equation}\label{eq:logconcave-functional}
 V(r)=\sum_{j=1}^d\int_{e_j^\perp}
                     \operatorname*{ess\,sup}_t r(z+te_j)\,dz,
\end{equation}
since an integrable unimodal slice has total variation twice its essential maximum. Boundary jumps of a compactly supported density are included.

Let $u,v$ be finite convex functions on $D$ with $u-v\in L^2(r)$ and gradients in $L^2(r)$. Put
\[
 \delta=\inf_{a\in\R}\norm{u-v-a}_{L^2(r)},\qquad
 X=\norm{\nabla u-\nabla v}_{L^2(r)},
\]
\[
 \tau_{u,v}(L)=\int_{\{\norm{\nabla u}>L\}}\norm{\nabla u}^2r
             +\int_{\{\norm{\nabla v}>L\}}\norm{\nabla v}^2r.
\]
All unlabelled integrals are with respect to Lebesgue measure on $D$.

\begin{theorem}[Tail-sensitive gradient interpolation]\label{thm:tail}
For every $h,L>0$,
\begin{equation}\label{eq:main-interpolation}
 \boxed{X^2\le12d\,\delta^2h^{-2}+28V(r)L^2h+4\tau_{u,v}(L).}
\end{equation}
Consequently, with
\[
 A_r=3\,2^{-2/3}(12d)^{1/3}(28V(r))^{2/3},
\]
one has
\begin{equation}\label{eq:infimum}
 \boxed{X^2\le\inf_{L>0}\{A_r\delta^{2/3}L^{4/3}+4\tau_{u,v}(L)\}.}
\end{equation}
When $\delta=0$, both conclusions hold with $X=0$.
\end{theorem}

The statement applies to arbitrary densities satisfying~\eqref{eq:slice-variation}; log-concavity is not needed for this theorem. Convexity of the potentials, rather than convexity of every density level set, is what makes the derivative truncation possible.

\begin{lemma}[One interval]\label{lem:interval}
Let $u,v$ be convex and $L$-Lipschitz on a bounded interval $I$, and let $f=u-v-a$. Then for every $h>0$,
\[
 \int_I|u'-v'|^2\le12h^{-2}\int_I f^2+28L^2h.
\]
The same assertion holds for open or half-open intervals and for restrictions whose endpoints are obtained by limits.
\end{lemma}
\begin{proof}
On $I_h=\{t\in I:t+h\in I\}$ write
$D_hu(t)=(u(t+h)-u(t))/h$. Convexity gives, almost everywhere,
\[
 0\le D_hu-u'\le u'(t+h)-u'(t)\le2L.
\]
Integrating the derivative difference and cancelling translates yields at most $2Lh$, so
$\int_{I_h}(D_hu-u')^2\le4L^2h$; likewise for $v$. Also
$\int_{I_h}|D_hf|^2\le4h^{-2}\int_I f^2$. Using
$u'-v'=D_hf-(D_hu-u')+(D_hv-v')$ gives $12h^{-2}\int_I f^2+24L^2h$ on $I_h$. The remaining interval length is at most $h$, and its squared derivative difference is at most $4L^2$. This proves the estimate, including $h\ge|I|$.
\end{proof}

\begin{proof}[Proof of Theorem~\ref{thm:tail}]
We first explain the one-dimensional level-set counting used below. For a nonnegative integrable $BV$ function $w$ on the line, almost every superlevel set agrees up to null sets with a finite union of bounded intervals. If $n_w(s)$ is the number of these intervals, then
\begin{equation}\label{eq:coarea}
 \int_0^\infty n_w(s)\,ds=\tfrac12\TV(w).
\end{equation}
This is the one-dimensional coarea formula~\cite{FR,Fleming}. It can also be seen directly by using
$|b-c|=\int_0^\infty|\ind_{b>s}-\ind_{c>s}|\,ds$ on nested refining partitions in the definition of variation. Monotone convergence gives the integral of the variation of the superlevel indicator; each bounded component contributes two jumps. Integrability rules out an unbounded component at a positive level.

Fix a coordinate line and a positive level $s$ for its density. On each interval component $I$ of that superlevel set, the derivatives of $u$ and $v$ are monotone where they exist. Their common good set
\[
 J=\{t\in I:|u'(t)|\le L,\ |v'(t)|\le L\}
\]
agrees up to a null set with an interval, possibly empty. On $J$ both restricted functions are $L$-Lipschitz. Lemma~\ref{lem:interval} therefore gives
\[
 \int_J|u'-v'|^2\le12h^{-2}\int_J|u-v-a|^2+28L^2h.
\]
On its complement, at least one derivative has magnitude exceeding $L$, and the pointwise bound
\[
 |u'-v'|^2\le4\bigl(|u'|^2\ind_{|u'|>L}
                       +|v'|^2\ind_{|v'|>L}\bigr)
\]
holds: if just one is large it dominates the magnitude of the other; if both are large use $|b-c|^2\le2(b^2+c^2)$. Thus each original component costs at most one boundary term. This is why clipping the derivatives does not create an uncontrolled number of intervals.

Sum over components, integrate over levels and over transverse lines, and use layer cake, Tonelli, and~\eqref{eq:coarea}. The $j$th coordinate contributes
\[
 \int|\partial_j u-\partial_j v|^2r
 \le12h^{-2}\int|u-v-a|^2r+28L^2hV_j(r)
 +4\int\bigl(|\partial_j u|^2\ind_{|\partial_j u|>L}
              +|\partial_j v|^2\ind_{|\partial_j v|>L}\bigr)r,
\]
where $\sum_j V_j(r)=V(r)$. Convex functions are locally Lipschitz in $D$ and differentiable almost everywhere, so the slice derivatives agree with the corresponding gradient components on a common full-measure set. Endpoints on $\partial D$ do not change the integrals. Summing in $j$, the component-tail sum is bounded by the vector tail $\tau_{u,v}(L)$. Infimize in $a$ to obtain~\eqref{eq:main-interpolation}.

For $\delta>0$, minimize the first two terms in $h$; the identity
$\min_{h>0}(Ah^{-2}+Bh)=3\,2^{-2/3}A^{1/3}B^{2/3}$ yields~\eqref{eq:infimum}. If $\delta=0$, take $h\downarrow0$ in~\eqref{eq:main-interpolation}, then $L\uparrow\infty$. The $L^2$ gradient assumption gives $\tau(L)\to0$, hence $X=0$. The nonnegative density has $V(r)>0$; otherwise all coordinate slices would be constant and integrability would force zero mass.
\end{proof}

\section{Finite moments and a sharp interpolation exponent}
\begin{theorem}\label{thm:moment}
Suppose, in addition, that for some $p>2$
\[
 M=\int(\norm{\nabla u}^p+\norm{\nabla v}^p)r<\infty.
\]
Set $\alpha_p=(p-2)/(3p-2)$. Then
\begin{equation}\label{eq:moment}
 X\le\sqrt2\,A_r^{\frac{3(p-2)}{2(3p-2)}}
                 (4M)^{\frac2{3p-2}}\delta^{\alpha_p}.
\end{equation}
The exponent $\alpha_p$ cannot be replaced by a larger exponent in a uniform inequality of this kind with a fixed source density and bounded $M$.
\end{theorem}
\begin{proof}
The tail is bounded by $\tau(L)\le ML^{2-p}$. For $M\delta>0$, balance the two terms in~\eqref{eq:infimum} by taking
\[
 L=\left(\frac{4M}{A_r\delta^{2/3}}\right)^{1/(p-2/3)}.
\]
Each term then equals
$(A_r\delta^{2/3})^{3(p-2)/(3p-2)}(4M)^{4/(3p-2)}$.
Taking the square root proves~\eqref{eq:moment}; zero cases follow by the previous theorem.

For sharpness take $r=\ind_{(0,1)}$, $v=0$, and
\[
 u_a(x)=a^{-1/p}(x-1+a)_+,\quad0<a<1.
\]
These are restrictions of globally finite convex functions. Direct integration gives
\[
 M=1,\quad X=a^{1/2-1/p},\quad
 \delta^2=a^{3-2/p}\left(\frac13-\frac a4\right).
\]
Thus $X/\delta^\alpha$ diverges for any
$\alpha>(1/2-1/p)/(3/2-1/p)=\alpha_p$. At $p=2$ the same example has $X=1$ while $\delta\to0$, showing that a bounded second gradient moment alone does not supply a uniform potential-to-gradient modulus. This last statement is about arbitrary convex gradients, not about stability in Wasserstein distance.
\end{proof}

\begin{remark}[Sharpness is not transferred automatically to transport]
In the one-dimensional example, one induced target is a point mass and the common-source optimal coupling gives $W_2=X$. It is therefore not a counterexample to a better Brenier-map exponent in $W_2$. The ramp proves sharpness only for interpolation. The different two-atom construction in Section~\ref{sec:steep} establishes transport sharpness for a fixed smooth source. A smooth positive density on $[0,1]$, including a truncated Gaussian, gives the same interpolation exponents after normalizing the ramp height by its $p$th moment, since its mass and weighted polynomial integrals near the endpoint have the corresponding leading powers of $a$.
\end{remark}

\section{Application to unbounded-target optimal transport}\label{sec:transport}
For an absolutely continuous probability source $\rho=r\,dx$, write $T_\mu=\nabla U_\mu$ for its quadratic Brenier map and centered convex potential. We separate the transport input from the interpolation theorem:
\begin{equation}\tag{P}\label{eq:P}
 \norm{U_\mu-U_\nu}_{L^2(\rho)}
 \le K_\rho W_2(\mu,\nu)\qquad(\mu,\nu\in\Ptwo(\R^d)).
\end{equation}
The centered representatives are assumed finite on the open convex source domain and in $L^2(\rho)$.

\begin{theorem}[Tail modulus for transport]\label{thm:transport}
Assume~\eqref{eq:slice-variation} and~\eqref{eq:P}. Write $w=W_2(\mu,\nu)$ and
\[
 \tau_{\mu,\nu}(L)=\int_{|y|>L}|y|^2\,d\mu(y)
                    +\int_{|y|>L}|y|^2\,d\nu(y).
\]
Then
\begin{equation}\label{eq:transport-tail}
 \boxed{\norm{T_\mu-T_\nu}_{L^2(\rho)}^2
 \le\inf_{L>0}\{A_rK_\rho^{2/3}w^{2/3}L^{4/3}
                         +4\tau_{\mu,\nu}(L)\}.}
\end{equation}
If $M_p=(\int|y|^p\,d\mu+\int|y|^p\,d\nu)^{1/p}<\infty$ for $p>2$, then
\begin{equation}\label{eq:transport-moment}
 \boxed{\norm{T_\mu-T_\nu}_{L^2(\rho)}
 \le C_{p,r,K_\rho,d}\,M_p^{\frac{2p}{3p-2}}
                         w^{\frac{p-2}{3p-2}}.}
\end{equation}
One may take
$C_{p,r,K_\rho,d}=\sqrt2 A_r^{3(p-2)/(2(3p-2))}4^{2/(3p-2)}K_\rho^{\alpha_p}$.
\end{theorem}
\begin{proof}
The pushforward identity $(T_\mu)_\#\rho=\mu$ identifies gradient tails and moments with target tails and moments. Apply Theorems~\ref{thm:tail} and~\ref{thm:moment}, then~\eqref{eq:P}. No target truncation is used in this step, and no unbounded Lipschitz constant is inserted into a bounded-target estimate.
\end{proof}

\begin{corollary}[Sources supplied by manuscript 001]\label{cor:sources}
The conclusions of Theorem~\ref{thm:transport} hold in either of the following cases:
\begin{enumerate}
\item $r=Z^{-1}e^{-\kappa|x|^2/2-W(x)}$, where $\kappa>0$ and $W$ is proper lower-semicontinuous convex with full-dimensional effective domain. Take $K_\rho=\sqrt{2/\kappa}$.
\item $\rho$ is full-dimensional log-concave, supported on a compact convex set in $B_R$. Take $K_\rho=(1+\sqrt{162})R$.
\end{enumerate}
No common bounded support of the targets, lower atom-mass bound, target density, or moment assumption $p>d$ is required.
\end{corollary}
\begin{proof}
In case (1), manuscript 001 v2, Theorem 1.1, supplies~\eqref{eq:P} for all finite-second-moment targets. Its density-envelope argument gives finite~\eqref{eq:logconcave-functional}. In case (2), its compact companion theorem supplies the same estimate uniformly over bounded targets with a constant independent of the target radius. The following completion argument extends that estimate to $\Ptwo$.

Project a target $\mu$ radially onto balls of radius $n$, obtaining $\mu_n\to\mu$ in $W_2$. The centered convex potentials are Cauchy in $L^2(\rho)$ by the bounded-target estimate. On every compact subset of the interior source domain, the density is positive and bounded below, so they converge in $L^1_{\rm loc}$. The elementary compactness theorem for finite convex functions then gives a finite convex representative of the limit, local uniform convergence, and almost-everywhere gradient convergence. For completeness, local $L^1$ control bounds a convex function above by its ball average; a reflection about a nearby point with bounded value bounds it below on a smaller ball. The resulting local oscillation bound bounds subgradients on a further smaller ball. Arzel\`a--Ascoli and the supporting-plane inequality give the stated conclusions.

Testing gradients against bounded continuous functions identifies their limiting law as $\mu$. The gradient is in $L^2(\rho)$ because that law has finite second moment. Extend the finite convex representative by its supporting affine functions. Fenchel equality on its gradient gives
$U^*(\nabla U)=x\cdot\nabla U-U$, which is integrable in absolute value under the source, using its second moment and $U\in L^2$. Fenchel's inequality is therefore integrable against any competing coupling and proves quadratic optimality of the gradient plan. Almost-everywhere differentiability gives uniqueness. Potentials with the same gradient differ by a constant on the connected source interior. Performing the construction for two targets and passing the bounded-target inequality to the limit establishes~\eqref{eq:P}. Compact log-concavity also gives finite slice functional, as in 001. This completion uses the source-dependent potential estimate; it is not an independent proof of that estimate.
\end{proof}

\begin{corollary}[Uniform integrability and linearized transport]
On any class $\mathcal C\subset\Ptwo(\R^d)$ with
\[
 \lim_{L\to\infty}\sup_{\mu\in\mathcal C}
                  \int_{|y|>L}|y|^2\,d\mu(y)=0,
\]
Theorem~\ref{thm:transport} gives a uniform map-continuity modulus through~\eqref{eq:transport-tail}. Moreover
\[
 W_2(\mu,\nu)\le\norm{T_\mu-T_\nu}_{L^2(\rho)},
\]
so on bounded $p$-moment classes, $p>2$, the common-source linearized transport distance is bi-H\"older equivalent to $W_2$, with upper exponent $\alpha_p$.
\end{corollary}
\begin{proof}
Choose $L$ to control both tails uniformly, and then choose $w$ small to control the first term. For the lower bound, $(T_\mu,T_\nu)_\#\rho$ is a coupling of the two targets. Its cost bounds the optimal coupling cost from above.
\end{proof}

\section{Stretched-exponential tails}
\begin{corollary}\label{cor:exponential}
Under the assumptions of Theorem~\ref{thm:transport}, suppose for some $q,K>0$ and $B<\infty$ that
\[
 \int e^{(|y|/K)^q}\,d\mu(y)+\int e^{(|y|/K)^q}\,d\nu(y)\le B.
\]
For $0<w=W_2(\mu,\nu)\le K$,
\[
 \norm{T_\mu-T_\nu}_{L^2(\rho)}
 \le C_{q,B,r,K_\rho,d}\,K^{2/3}w^{1/3}
                      \bigl(1+\log(K/w)\bigr)^{2/(3q)}.
\]
In particular, a common subgaussian target-tail bound gives a logarithmic factor with exponent $1/3$.
\end{corollary}
\begin{proof}
The finite constant $c_q=\sup_{s\ge0}s^2e^{-s^q/2}$ gives
\[
 \tau_{\mu,\nu}(L)\le c_q B K^2 e^{-(L/K)^q/2}.
\]
Indeed, for $s>L/K$, $s^2\le c_qe^{s^q/2}\le c_qe^{s^q}e^{-(L/K)^q/2}$, and the assumed exponential moment can be integrated. Put $z=w/K\in(0,1]$ and take
\[
 L=K\bigl(2+\tfrac43\log(1/z)\bigr)^{1/q}.
\]
Then the tail is at most $c_qB e^{-1}K^2z^{2/3}$. The first term of~\eqref{eq:transport-tail} is
$A_rK_\rho^{2/3}K^2z^{2/3}(2+\tfrac43\log(1/z))^{4/(3q)}$.
Taking square roots proves the assertion. No claim is made that the logarithmic factor is sharp.
\end{proof}

\section{Translation regularity beyond bounded Hessians}\label{sec:translation}
The parallel manuscript 001 v3~\cite{M001v3} proves an unclipped Gaussian interpolation inequality and extends it to full-support strongly convex $C^{1,1}$ source potentials with globally bounded Hessian. Its signed-square and minimum-weight argument suggests the following explicit hypothesis. We reproduce that argument to specify precisely what is used, then verify the hypothesis for superquadratic sources not covered by a global Hessian bound.

Let $r>0$ on $\R^d$, let $q>2$ and $s=q/(q-2)$, and suppose
\begin{equation}\label{eq:translation-assumption}
 \mathcal R_s(r):=\max_{1\le j\le d}\sup_{0<h\le1,\,\epsilon\in\{-1,1\}}
 \frac1h\left\|\frac{r(\cdot+\epsilon he_j)}r-1\right\|_{L^s(r)}<\infty.
\end{equation}
This is an integrable translation-ratio condition, not merely pointwise smoothness.

\begin{proposition}[Translation-ratio interpolation]\label{prop:translation}
For finite convex $u,v$ on $\R^d$, gradients in $L^q(r)$, and $u-v-a\in L^2(r)$, condition~\eqref{eq:translation-assumption} implies, for $0<h\le1$,
\[
 \norm{\nabla u-\nabla v}_{L^2(r)}^2
 \le12d h^{-2}\norm{u-v-a}_{L^2(r)}^2
 +8h\mathcal R_s(r)\sum_{j=1}^d
 \bigl(\norm{\partial_j u}_{L^q(r)}^2+
       \norm{\partial_j v}_{L^q(r)}^2\bigr).
\]
If the source also satisfies the potential estimate~\eqref{eq:P}, then, with
$M=\norm{T_\mu}_{L^q(r)}+\norm{T_\nu}_{L^q(r)}$,
\begin{equation}\label{eq:translation-map}
 \norm{T_\mu-T_\nu}_{L^2(r)}
 \le\sqrt{d(12K_\rho^2+8\mathcal R_s(r))}\,
             M^{2/3}W_2(\mu,\nu)^{1/3}.
\end{equation}
\end{proposition}
\begin{proof}
Fix a coordinate and set $m_h(x)=\min(r(x),r(x+he_j))$.
Writing $f=u-v-a$, both endpoints of its difference quotient are controlled by their original weights, so
\[
 \int|D_{j,h}f|^2m_h\le4h^{-2}\norm f_{L^2(r)}^2.
\]
Let $z=\partial_j u$. Convexity gives
$0\le D_{j,h}u-z\le z(x+he_j)-z(x)$.
For real $A\le B$,
$(B-A)^2\le2(B|B|-A|A|)$.
With $g=z|z|$, monotonicity and $m_h\le r$ therefore give
\[
 \int|D_{j,h}u-z|^2m_h
 \le2\int(g(x+he_j)-g(x))r(x)\,dx
 \le2h\mathcal R_s(r)\norm z_{L^q(r)}^2.
\]
For the last step change variables in the first term, then apply H\"older with exponents $q/2,s$. Absolute integrability of the translated signed-square functions follows from the same hypothesis: their integrals are bounded by the original integral plus
$h\mathcal R_s\norm g_{L^{q/2}(r)}$.
The same calculation applies to $v$. The three-term square inequality for the difference quotient and its two errors gives coefficient six for each derivative norm on $m_h$. On the remaining weight,
\[
 \int|\partial_j u-\partial_j v|^2(r-m_h)
 \le2h\mathcal R_s(r)
       (\norm{\partial_j u}_q^2+\norm{\partial_j v}_q^2).
\]
This proves coefficient eight and the asserted interpolation after summing coordinates.

For transport, the coordinate norm sum is at most $dM^2$. Also $w=W_2(\mu,\nu)\le M$, by the triangle inequality through $\delta_0$ and the $L^q$ moment bound. If $w>0$, put $h=(w/M)^{2/3}\le1$ and use~\eqref{eq:P}; this gives~\eqref{eq:translation-map}. Zero cases follow from uniqueness.
\end{proof}

\begin{corollary}[Superquadratic confinement]\label{cor:polynomial}
Let $r=Z^{-1}e^{-V}$, where $V\in C^1(\R^d)$ is $\kappa$-strongly convex. Suppose for some $m>1$ and finite positive constants $c,C$,
\[
 V(x)\ge c|x|^m-C,\qquad
 |\nabla V(x)|\le C(1+|x|)^{m-1}.
\]
Then for every $q>2$, $\mathcal R_{q/(q-2)}(r)<\infty$, and the one-third estimate~\eqref{eq:translation-map} holds with $K_\rho=\sqrt{2/\kappa}$.
In particular this covers $V(x)=\kappa|x|^2/2+|x|^m$ for every real $m>2$, although its Hessian is unbounded.
\end{corollary}
\begin{proof}
For either unit coordinate shift of size $h\le1$, the fundamental theorem of calculus gives
$|V(x\pm he_j)-V(x)|\le hG(x)$ with
$G(x)=C(2+|x|)^{m-1}$. Hence
\[
 \left|r(x\pm he_j)/r(x)-1\right|
 \le hG(x)e^{G(x)}.
\]
For every finite $s$, $(G e^G)^s$ is integrable against $r$: its exponential growth is $O(|x|^{m-1})$, strictly smaller than the confining order $c|x|^m$, and its polynomial factor is harmless. This proves~\eqref{eq:translation-assumption}. Strong log-concavity supplies~\eqref{eq:P} from 001. The displayed example is strongly convex and has the stated growth bounds.
\end{proof}

\begin{corollary}[Product superquadratic sources]\label{cor:product}
For every fixed source
\[
 r(x)=\prod_{j=1}^d Z_j^{-1}
       \exp(-\kappa_jx_j^2/2-\lambda_j|x_j|^{m_j}),
 \quad \kappa_j,\lambda_j>0,\quad m_j>2,
\]
the map estimate has exponent $1/3$ for every finite target moment order $q>2$. If $d\ge2$, that exponent is sharp for each such fixed source, even for three-atom targets in $B_1$.
\end{corollary}
\begin{proof}
A coordinate translation changes only its corresponding marginal. The preceding growth argument applies separately with exponent $m_j$, proving the translation-ratio condition despite possibly unequal growth orders. The source is strongly log-concave with parameter $\min_j\kappa_j$.

For sharpness let $X,Z$ be its first two independent coordinates and put $F(t)=\Pr(|X|<t)$ for $t\ge0$. The even density of $X$ is bounded and has value $f_1(0)>0$, so $F(t)/t\to2f_1(0)$ and $F(t)\le C t$. Fix $k>0$, set $b=ka$, and choose $c_a$ from
\[
 \E F((c_a+kaZ)_+)=F(a)=:m_a.
\]
The left side increases continuously from zero to one, and has a unique solution at the positive level $F(a)$. Dominated convergence and monotone bracketing give $c_a/a\to c_0$, where $\E(c_0+kZ)_+=1$. The finite first moment of $Z$ suffices here. Its nondegeneracy implies
\[
 \frac1a\E|F((c_a+kaZ)_+)-F(a)|
 \longrightarrow2f_1(0)\E|(c_0+kZ)_+-1|>0.
\]
Use convex potentials $u_a=\max(-x_1,x_1,a)$ and
$v_a=\max(-x_1,x_1,c_a+ka x_2)$. Symmetry in $x_1$ gives the same outer masses. Their targets have outer atoms $\pm e_1$ and central atoms $0$ and $be_2$, respectively, all in $B_1$ for small $a$. Exactly as in the three-atom construction of 001, the target distance satisfies $W_2^2=m_ab^2\asymp a^3$, whereas the squared map discrepancy is the central-cell symmetric difference plus $m_ab^2$, and is $\asymp a$. This excludes every exponent above $1/3$. The calculation identifies every needed property of the new source; it does not appeal to class-level sharpness via a different Gaussian source.
\end{proof}

\section{A smooth full-support source with boundary-type instability}\label{sec:steep}
The previous positive result cannot be deduced from smoothness, full support and strong log-concavity alone. We construct one fixed smooth source for which the general finite-moment exponent is optimal. The source does not depend on the moment order.

\begin{theorem}[Smoothness and full support are insufficient]\label{thm:steep}
For every $d\ge2$ there is a fixed strictly positive $C^\infty$ probability density $r$ on $\R^d$, with $-\log r$ having Hessian at least the identity, such that for every $q>2$ and every
$\alpha>(q-2)/(3q-2)$,
\[
 \sup_{\substack{\mu,\nu\text{ two-atom targets}\\
       \int|y|^q\,d\mu=\int|y|^q\,d\nu=1}}
 \frac{\norm{T_\mu-T_\nu}_{L^2(r)}}{W_2(\mu,\nu)^\alpha}
 =\infty.
\]
Together with Theorem~\ref{thm:transport}, this makes $(q-2)/(3q-2)$ the optimal power for each fixed $q>2$ for this same source. For $q=2$, the construction gives exact unit-second-moment targets with $W_2\to0$ but map discrepancy bounded away from zero. No boundary of the source support is present.
\end{theorem}

\subsection*{An explicit locally finite sequence of smooth steep layers}
Let $\theta\in C^\infty(\R)$ be the increasing step equal to zero on $(-\infty,0]$, one on $[1,\infty)$, and on $(0,1)$ given by
\[
 \theta(t)=\frac{e^{-1/t}}{e^{-1/t}+e^{-1/(1-t)}}.
\]
All its derivatives are flat at the endpoints, and $\theta(1-t)=1-\theta(t)$. Put $\Psi(t)=\int_0^t\theta(s)\,ds$. Then $\Psi$ is smooth and convex, is zero for $t\le0$, and equals $t-1/2$ for $t\ge1$.

Define recursively, for integers $n\ge2$,
\[
 B_n=\frac{n^2}{2}+\sum_{2\le j<n}M_j(n-j-\varepsilon_j/2),
 \qquad D_n=n+\sum_{2\le j<n}M_j,
\]
\begin{equation}\label{eq:layer-parameters}
 N_n=n^2\bigl(1+\lceil B_n+D_n\rceil\bigr),\quad
 a_n=2^{-N_n},\quad M_n=a_n^{-2},\quad\varepsilon_n=a_n^4.
\end{equation}
Every step is an exact rational or integer definition, although the numbers grow extremely rapidly. Define
\begin{equation}\label{eq:steep-potential}
 V(x)=x^2/2+\sum_{n\ge2}M_n\varepsilon_n
                       \Psi((x-n)/\varepsilon_n),\qquad
 r_1(x)=Z_1^{-1}e^{-V(x)}.
\end{equation}
On every compact set only finitely many summands are nonzero. Hence $V$ is finite and $C^\infty$, with $V''\ge1$. All added terms are nonnegative, so $0<Z_1\le\sqrt{2\pi}$; a positive lower bound follows by integration over $(-1,0)$. Moreover $V(n)=B_n$ and $V'(n)=D_n$. Since every $\varepsilon_j<1/4$, all earlier layers are in their linear regimes at $n$. The recursion implies
\begin{equation}\label{eq:layer-limits}
 a_nD_n\longrightarrow0,\quad
 \frac{B_n}{|\log a_n|}\longrightarrow0,\quad
 \frac{\varepsilon_n}{a_n}\longrightarrow0,\quad
 \frac1{M_na_n}\longrightarrow0.
\end{equation}
Indeed $N_n\ge n^2(1+B_n+D_n)$ and $|\log a_n|=N_n\log2$; the last two ratios are $a_n^3$ and $a_n$.

\subsection*{Local tails converge to a hard-edge profile without a hard edge}
Let $X$ have density $r_1$ and put $\beta_n=r_1(n)$. Define
\[
 G_n(t)=\frac{\Pr(X>n-a_nt)}{\beta_na_n}\quad(t\in\R).
\]
For every fixed real $t$,
\begin{equation}\label{eq:profile-limit}
 G_n(t)\longrightarrow t_+,
 \qquad
 0\le G_n(t)\le t_+e^{t_+}+1.
\end{equation}
To prove this, the mass to the right of $n$ is bounded by
\[
 \Pr(X>n)\le\beta_n(\varepsilon_n+M_n^{-1}):
\]
the layer interval costs at most $\beta_n\varepsilon_n$, and afterward $V'\ge M_n$, so its remaining tail costs at most $\beta_n/M_n$. After division by $\beta_na_n$ this tends to zero. For $t>0$ the remaining interval is $[n-a_nt,n]$. For fixed $t$ it lies between consecutive layers for all large $n$, where
\[
 V(n-a_ns)=B_n-D_na_ns+a_n^2s^2/2\quad(0\le s\le t).
\]
Equation~\eqref{eq:layer-limits} makes its normalized integral tend to $t$. For arbitrary $t>0$, convexity at $n$ gives
$r_1(n-a_ns)/\beta_n\le e^{D_na_ns}$. Since $D_na_n\le1$, the global bound in~\eqref{eq:profile-limit} follows, including the vanishing right tail; for $t\le0$ use that right-tail bound alone.

Set the $d$-dimensional source to be
\[
 r(x)=r_1(x_1)\prod_{j=2}^d\varphi(x_j),
 \qquad \varphi(t)=(2\pi)^{-1/2}e^{-t^2/2}.
\]
It has the smoothness, positivity and convexity properties in the theorem. Its slice functional is finite: $V(r)=Z_1^{-1}+(d-1)/\sqrt{2\pi}$, since $V$ has its minimum zero at the origin and each marginal is unimodal.

\subsection*{Two rare atoms with exactly normalized moments}
Fix $k>0$, let $Z\sim N(0,1)$ be independent of $X$, and set
$m_n=\Pr(X>n-a_n)=\beta_na_nG_n(1)$ and $b_n=ka_n$.
Choose $s_n\in\R$ to match the half-space masses:
\begin{equation}\label{eq:layer-mass-match}
 \E G_n(s_n+kZ)=G_n(1).
\end{equation}
The unnormalized tail on the left is continuous and strictly increasing from zero to one as $s_n$ increases, so the solution exists uniquely. The majorant in~\eqref{eq:profile-limit} is Gaussian-integrable for $t=s+kZ$, uniformly for bounded $s$. Dominated convergence and monotone bracketing yield
\[
 s_n\to s_0,\qquad \E(s_0+kZ)_+=1.
\]
In particular $(s_n)$ is bounded. For the cells
\[
 A_n=\{x_1>n-a_n\},\qquad
 C_n=\{x_1+b_nx_2>n-a_ns_n\},
\]
both masses equal $m_n$, and the nested intervals on each $x_2$ slice give
\begin{equation}\label{eq:layer-symdiff}
 \frac{r(A_n\mathbin\triangle C_n)}{m_n}
 =\frac{\E|G_n(s_n+kZ)-G_n(1)|}{G_n(1)}
 \longrightarrow d_0:=\E|(s_0+kZ)_+-1|>0.
\end{equation}
This limit uses the same Gaussian-integrable majorant. Nondegeneracy of $Z$ makes $d_0$ strictly positive.

For $q\ge2$, put $R_n=m_n^{-1/q}$ and
$v_n=(e_1+b_ne_2)/\sqrt{1+b_n^2}$. Use targets
\[
 \mu_n=(1-m_n)\delta_0+m_n\delta_{R_ne_1},\qquad
 \nu_n=(1-m_n)\delta_0+m_n\delta_{R_nv_n}.
\]
Their $q$th moments are exactly one. Their convex potentials are
\[
 u_n(x)=R_n(x_1-n+a_n)_+,\qquad
 w_n(x)=\frac{R_n}{\sqrt{1+b_n^2}}
             (x_1+b_nx_2-n+a_ns_n)_+.
\]
Their gradients have the stated target laws by~\eqref{eq:layer-mass-match}, and integrating convex supporting inequalities proves quadratic optimality. All functions have at most linear growth for each fixed $n$, so source integrability causes no difficulty.

Let $\Delta_n=|e_1-v_n|\sim ka_n$. A target coupling sends some mass $t\le m_n$ between the two rare atoms and the remaining rare mass through zero. Its cost is
$R_n^2[t\Delta_n^2+2(m_n-t)]$. Since $\Delta_n^2<2$, the optimum is $t=m_n$, which is feasible. Hence
\[
 W_2(\mu_n,\nu_n)^2=m_nR_n^2\Delta_n^2.
\]
The exact map-cost identity is
\[
 \norm{T_{\mu_n}-T_{\nu_n}}_{L^2(r)}^2
 =R_n^2\bigl(r(A_n\mathbin\triangle C_n)
                  +\Delta_n^2r(A_n\cap C_n)\bigr).
\]
Writing $\vartheta_q=1/2-1/q$, equation~\eqref{eq:layer-symdiff} gives
\begin{equation}\label{eq:layer-costs}
 \norm{T_{\mu_n}-T_{\nu_n}}_2\sim\sqrt{d_0}\,m_n^{\vartheta_q},
 \qquad W_2(\mu_n,\nu_n)\sim k a_nm_n^{\vartheta_q}.
\end{equation}
Finally $m_n=Z_1^{-1}e^{-B_n}a_n(1+o(1))$, so~\eqref{eq:layer-limits} implies
\[
 \frac{\log m_n}{\log a_n}\longrightarrow1.
\]
For any $\alpha>\vartheta_q/(1+\vartheta_q)=(q-2)/(3q-2)$, the logarithm of the ratio in Theorem~\ref{thm:steep}, divided by $|\log a_n|$, tends to
$\alpha-\vartheta_q(1-\alpha)>0$. This proves divergence. At $q=2$,~\eqref{eq:layer-costs} gives map distance tending to $\sqrt{d_0}>0$ while $W_2\to0$. The source is the same for every $q$, and the matching upper exponent follows from Theorem~\ref{thm:transport} and the strong-log-concavity potential estimate. This completes the proof of Theorem~\ref{thm:steep}.

\begin{remark}[What the counterexample distinguishes]
The loss from $1/3$ is not a phenomenon exclusive to hard source boundaries. The source above is smooth and has full support, but its derivatives contain increasingly steep layers and its Hessian is unbounded. In contrast, polynomial confinement and the globally bounded-Hessian class satisfy an integrable translation-ratio bound and retain $1/3$. The construction is a separate two-atom example, not an inference from the one-dimensional interpolation ramp. Its lower bounds require no unproved source stability input; that input is used only for the matching upper estimate.
\end{remark}

\section{Scope of verification and of the claimed advance}
The exact verifier accompanying this manuscript checks the boundary-ramp identities, the exponent-balancing identities, and one-dimensional piecewise-affine instances of~\eqref{eq:main-interpolation}, with rational slopes, breakpoints, thresholds, and step sizes. It uses integer and rational arithmetic only. These tests do not formalize the coarea argument or the potential estimate imported from 001.

The sharp statement is the potential-to-gradient interpolation exponent under a $p$th moment bound. The finite-moment \emph{transport} exponent is sharp for the fixed smooth source of Theorem~\ref{thm:steep}, not for every source; the product sources in Corollary~\ref{cor:product} have the better sharp exponent $1/3$. The second-moment interpolation counterexample must not be presented as a failure of uniform Brenier stability on a bounded $W_2$ set. For comparison, Delalande--M\'erigot~\cite{DM} treats unbounded targets under moment conditions including $p>d$ with a compact source bounded above and below; our corollary changes both the source hypotheses and the moment threshold and uses a different potential estimate. A comprehensive comparison with all later papers has not been completed.

\clearpage
\begin{thebibliography}{9}
\bibitem{M001} mxym, \emph{Sharp one-third stability of Brenier maps for strongly log-concave sources}, v2, October 7, 2026, manuscript 001, \texttt{mxym/math}. Audited source blob: \texttt{d6896f8ac49e3a945324d5d57c54e1797775b1e5}; pinned repository snapshot \texttt{5c6c088aa5abf1c1a4bdca6a8ce5beaa27faaef3}. AI-assisted research draft.
\bibitem{M001v3} mxym, \emph{Sharp Brenier stability under target moment bounds}, v3, October 7, 2026, manuscript 001. Pinned commit \texttt{5c6c088aa5abf1c1a4bdca6a8ce5beaa27faaef3}; source blob \texttt{0007edf5e28e4b4b0d34b2043cffd6e13e9f5aed}. This parallel version precedes the full source publication of the present v1 and contains the overlapping tail and Gaussian results acknowledged above.
\bibitem{O374} OpenAI, \emph{Sharp One-Third Stability of Brenier Maps}, family 374, \texttt{openai/math}, pinned snapshot \texttt{adc7f1241b42e322a6451854ab7e4b4c146bf78a}. The required finite-difference argument is reproduced above; this item is attributed through manuscript 001.
\bibitem{DM} A. Delalande and Q. M\'erigot, \emph{Quantitative stability of optimal transport maps under variations of the target measure}, Duke Mathematical Journal (2023), DOI: \href{https://doi.org/10.1215/00127094-2022-0106}{10.1215/00127094-2022-0106}; arXiv:2103.05934.
\bibitem{FR} W. H. Fleming and R. Rishel, \emph{An integral formula for total gradient variation}, Archiv der Mathematik 11 (1960), 218--222.
\bibitem{Fleming} W. H. Fleming, \emph{Functions whose partial derivatives are measures}, Illinois Journal of Mathematics 4(3) (1960), 452--478, DOI: \href{https://doi.org/10.1215/ijm/1255456061}{10.1215/ijm/1255456061}.
\end{thebibliography}
\end{document}
